\documentclass{article}


\usepackage{arxiv}

\usepackage[utf8]{inputenc} % allow utf-8 input
\usepackage[T1]{fontenc}    % use 8-bit T1 fonts
\usepackage{hyperref}       % hyperlinks
\usepackage{url}            % simple URL typesetting
\usepackage{booktabs}       % professional-quality tables
\usepackage{amsfonts}       % blackboard math symbols
\usepackage{nicefrac}       % compact symbols for 1/2, etc.
\usepackage{microtype}      % microtypography
\usepackage{lipsum}
\usepackage{graphicx}
\graphicspath{ {./images/} }

\title{Can We Trust LLM Self-Explanations? An Empirical Evaluation of Faithfulness in Open-Source Large Language Models}

\author{
Yash Sakharam Gavade\\
Department of Computer Science\\
Universität Trier\\
Trier, Germany\\
\texttt{s4yagava@uni-trier.de}
}

\begin{document}
\maketitle


\begin{abstract}
Large Language Models (LLMs) increasingly generate natural language explanations alongside their predictions, aiming to improve transparency and user trust. However, recent research has questioned whether these self-generated explanations faithfully represent the underlying reasoning process or merely provide plausible post-hoc justifications. This paper presents an empirical evaluation of the faithfulness of self-generated explanations produced by recent open-source LLMs. We compare multiple instruction-tuned models across diverse reasoning tasks using a unified prompting strategy and a lightweight evaluation protocol that assesses answer correctness, explanation faithfulness, consistency, and common explanation failure patterns. The study combines quantitative evaluation with qualitative error analysis to identify strengths and limitations of current self-explanation capabilities. The findings are expected to provide insights into the trustworthiness of open-source LLM explanations and offer practical recommendations for future explainable AI research.

\textbf{Note:} The abstract will be updated after completing the experimental evaluation to accurately summarize the final methodology, results, and contributions.

\end{abstract}



\keywords{
Explainable Artificial Intelligence,
Large Language Models,
Explainability,
Faithfulness,
Trustworthiness,
Open-Source LLMs,
Natural Language Processing,
Empirical Evaluation
}




\section{Introduction}

Large Language Models (LLMs) have fundamentally transformed the field of Natural Language Processing (NLP) by demonstrating remarkable capabilities across a wide range of language understanding and generation tasks. Recent open-source instruction-tuned models such as Llama, Qwen, Gemma, and DeepSeek have achieved competitive performance in question answering, reasoning, summarization, code generation, and conversational AI. Beyond generating accurate predictions, these models increasingly produce natural language explanations intended to justify or describe the reasoning behind their answers. Such explanations are commonly referred to as \emph{self-explanations}, as they are generated by the same model that produces the final prediction.

The ability of LLMs to explain their outputs has attracted significant attention within Explainable Artificial Intelligence (XAI). Explainability is considered an essential property for deploying AI systems in domains where transparency, accountability, and user trust are critical, including healthcare, education, finance, and legal decision support. In these settings, users often rely not only on the correctness of a prediction but also on the accompanying explanation when deciding whether to trust the model's output. Consequently, natural language explanations generated by LLMs are increasingly perceived as evidence of reasoning and transparency.

Despite their apparent usefulness, recent research has questioned whether self-generated explanations truly reflect the reasoning process used by LLMs to arrive at their predictions. Several studies have suggested that explanations may instead function as plausible post-hoc rationalizations, providing convincing justifications that do not necessarily correspond to the internal mechanisms responsible for generating the final answer. As a result, an explanation that appears logically coherent may still be unfaithful to the model's actual decision-making process. This distinction between plausibility and faithfulness has emerged as one of the central challenges in modern Explainable AI research.

Existing work has explored various aspects of explanation quality, including Chain-of-Thought reasoning, rationale generation, explanation consistency, and human evaluation of explanations. More recently, researchers have proposed intervention-based and causal evaluation methods to investigate explanation faithfulness. However, current literature remains fragmented. Many studies evaluate only a single model, focus exclusively on proprietary systems, or adopt different datasets, prompting strategies, and evaluation protocols, making direct comparison difficult. Furthermore, relatively few empirical studies systematically compare recent open-source LLMs under a unified evaluation framework.

To address this gap, this paper presents an empirical evaluation of the faithfulness of self-generated explanations produced by recent open-source Large Language Models. Specifically, multiple instruction-tuned LLMs are evaluated across diverse reasoning tasks using a standardized prompting strategy. The generated explanations are assessed using a lightweight evaluation protocol that combines quantitative performance measures with structured qualitative analysis. In addition to measuring answer correctness, the study investigates explanation faithfulness, consistency, and common explanation failure patterns observed across different reasoning tasks.

The main contributions of this work are summarized as follows:

\begin{itemize}
    \item We conduct a comparative empirical evaluation of self-generated explanations across multiple recent open-source Large Language Models.
    
    \item We introduce a lightweight evaluation protocol for assessing explanation faithfulness, consistency, and explanation quality within the scope of this study.
    
    \item We perform a qualitative error analysis to identify common explanation failure patterns across different reasoning tasks.
    
    \item We discuss the implications of explanation faithfulness for trustworthy and transparent deployment of open-source LLMs.
\end{itemize}

The remainder of this paper is organized as follows. Section~2 reviews related work on Explainable Artificial Intelligence, explainability in Large Language Models, and recent research on explanation faithfulness. Section~3 presents the proposed methodology, including the research design, selected models, datasets, prompting strategy, and evaluation framework. Section~4 describes the experimental setup. Section~5 presents the experimental results and qualitative analyses. Finally, Sections~6--9 discuss the findings, limitations, future research directions, and concluding remarks.




\section{Related Work}

\subsection{Explainable Artificial Intelligence}

Explainable Artificial Intelligence (XAI) aims to develop machine learning systems whose predictions can be understood, interpreted, and trusted by human users. As artificial intelligence models become increasingly complex, particularly with the emergence of deep neural networks, understanding the reasoning behind model predictions has become an important research challenge. Traditional machine learning models such as decision trees and linear regression provide relatively transparent decision-making processes, whereas modern deep learning architectures often function as black-box systems whose internal reasoning remains difficult to interpret.

Early XAI research primarily focused on post-hoc explanation techniques designed for computer vision and structured data. Popular approaches include Local Interpretable Model-Agnostic Explanations (LIME), SHapley Additive exPlanations (SHAP), Integrated Gradients, saliency maps, and attention visualization methods. These techniques attempt to explain model predictions by identifying influential input features or estimating feature importance without modifying the underlying predictive model.

While these approaches have proven valuable for conventional supervised learning models, their direct application to Large Language Models presents several limitations. Modern LLMs generate free-form natural language rather than fixed class labels, making traditional feature attribution techniques insufficient for explaining complex reasoning processes. Consequently, recent XAI research has shifted toward evaluating natural language explanations generated directly by LLMs themselves, giving rise to new research questions regarding explanation quality, interpretability, and faithfulness.

\subsection{Explainability in Large Language Models}

The rapid development of Large Language Models has significantly expanded the scope of explainability research. Unlike traditional classification models, LLMs are capable of generating detailed natural language explanations together with their predictions. These explanations often resemble human reasoning, making them attractive for improving transparency and user trust.

Several explanation strategies have emerged in recent years. Prompting techniques such as Chain-of-Thought (CoT) prompting encourage models to generate intermediate reasoning steps before producing the final answer. Other approaches include rationale generation, self-reflection, self-verification, debate-based reasoning, and tool-assisted reasoning. These methods aim to improve both reasoning performance and explanation quality while providing users with greater insight into the model's decision-making process.

Although explanation generation has become increasingly common in state-of-the-art LLMs, researchers have observed that explanation quality cannot be assessed solely by linguistic fluency or logical coherence. Explanations may appear convincing despite failing to accurately represent the reasoning process responsible for the final prediction. As a result, explainability research has gradually shifted from generating explanations toward evaluating whether those explanations are genuinely trustworthy.

\subsection{Evaluating Faithfulness of LLM Explanations}

Faithfulness has emerged as one of the central evaluation criteria for explanation quality in Large Language Models. An explanation is generally considered faithful if it accurately reflects the reasoning process that produced the model's prediction rather than merely providing a plausible justification after the answer has already been generated.

Recent studies have demonstrated that LLM-generated explanations may often function as post-hoc rationalizations. In many cases, models are capable of producing convincing explanations even when the corresponding prediction is incorrect or when the explanation does not align with the model's internal computation. This phenomenon raises concerns regarding user trust, particularly in high-stakes domains where explanations may influence important decisions.

To address this challenge, recent research has proposed several approaches for evaluating explanation faithfulness. These include intervention-based evaluation, causal analysis, perturbation methods, rationale consistency, explanation stability, and behavioral testing. Rather than evaluating explanations based solely on readability or human preference, these methods attempt to determine whether modifying or removing parts of the explanation changes the model's prediction, thereby providing stronger evidence of faithful reasoning.

Despite these advances, explanation faithfulness remains difficult to measure consistently across different LLMs. Existing studies often employ different datasets, prompting strategies, evaluation protocols, and proprietary models, making direct comparison challenging. Consequently, there is still a need for systematic empirical evaluations that compare recent open-source LLMs under identical experimental conditions.

\subsection{Research Gap and Motivation}

Although recent studies have substantially advanced the understanding of explanation generation and faithfulness evaluation in Large Language Models, several important research gaps remain. First, many existing investigations evaluate only a single LLM or focus primarily on proprietary commercial systems, limiting reproducibility and comparative analysis. Second, differences in prompting strategies, datasets, and evaluation methodologies make it difficult to directly compare findings across studies. Third, relatively little work has systematically investigated the behavior of multiple recent open-source LLMs using a unified experimental protocol specifically designed to evaluate self-generated explanations.

Furthermore, existing evaluation approaches frequently emphasize answer correctness while providing comparatively limited analysis of explanation consistency and common explanation failure patterns. Understanding why explanations fail is equally important for improving trustworthy AI systems, yet qualitative error analyses remain relatively scarce.

Motivated by these limitations, this work presents a unified empirical evaluation of explanation faithfulness across multiple recent open-source Large Language Models. By employing identical reasoning tasks, standardized prompts, and a consistent evaluation protocol, this study aims to provide a more comprehensive comparison of explanation quality while identifying common strengths and weaknesses of current self-explanation capabilities.





\section{Methodology}

\subsection{Research Design}

This study adopts a comparative empirical research design to evaluate the faithfulness of self-generated explanations produced by recent open-source Large Language Models (LLMs). Rather than proposing a new explanation generation method, the objective is to systematically investigate whether explanations generated by different LLMs faithfully justify their predicted answers across multiple reasoning tasks.

The overall methodology consists of four stages. First, a set of representative reasoning datasets covering different categories of reasoning is selected. Second, each dataset is evaluated using multiple open-source instruction-tuned LLMs under a standardized prompting strategy that requires the model to generate both an answer and a natural language explanation. Third, the generated outputs are evaluated using a unified evaluation protocol combining quantitative performance measures and structured qualitative assessment. Finally, the collected results are analyzed to compare explanation faithfulness across models and identify common explanation failure patterns.

A controlled experimental setting is adopted to ensure fair comparison between models. All evaluated LLMs receive identical prompts, identical input questions, and identical inference settings whenever possible. This standardized experimental design minimizes methodological variation and enables direct comparison of explanation quality across different models.

\subsection{Research Questions and Hypotheses}

The primary objective of this study is to investigate whether self-generated explanations produced by open-source LLMs faithfully represent the reasoning underlying their final predictions. To guide the empirical investigation, the following research questions are considered:

\begin{itemize}
    \item \textbf{RQ1:} How faithful are self-generated explanations produced by recent open-source Large Language Models?
    
    \item \textbf{RQ2:} Does explanation faithfulness differ across different open-source LLMs?
    
    \item \textbf{RQ3:} What explanation failure patterns are most commonly observed across different reasoning tasks?
\end{itemize}

Based on existing literature, the following hypotheses are formulated:

\begin{itemize}
    \item \textbf{H1:} Self-generated explanations are not always faithful to the reasoning process that produced the model's prediction.
    
    \item \textbf{H2:} Reasoning-oriented LLMs produce more faithful explanations than general instruction-tuned LLMs.
    
    \item \textbf{H3:} Explanation failure patterns vary across different categories of reasoning tasks.
\end{itemize}

These research questions and hypotheses form the basis of the experimental evaluation presented in the subsequent sections.

\subsection{Selected Open-Source Large Language Models}

This study focuses exclusively on recent open-source instruction-tuned Large Language Models to ensure reproducibility and accessibility. Open-source models allow complete transparency regarding model architecture, inference settings, and experimental procedures, making them particularly suitable for empirical explainability research.

Three representative LLMs from different model families are selected for evaluation. The selected models are expected to represent diverse architectural designs and reasoning capabilities while remaining computationally feasible for academic experimentation.

\begin{table}[h]
\centering
\caption{Selected Open-Source Large Language Models}
\begin{tabular}{lll}
\hline
Model & Organization & Parameters\\
\hline
Llama 3.2 Instruct & Meta & TODO\\
Qwen 2.5 Instruct & Alibaba & TODO\\
DeepSeek-R1 Distill & DeepSeek AI & TODO\\
\hline
\end{tabular}
\label{tab:models}
\end{table}

The final model versions and parameter sizes will be selected based on computational availability while ensuring comparable instruction-following capabilities.

\subsection{Reasoning Datasets}

To evaluate explanation faithfulness across different reasoning domains, this study employs publicly available benchmark datasets covering mathematical reasoning, commonsense reasoning, and multi-hop reasoning. Using multiple reasoning tasks allows investigation of whether explanation quality varies depending on the complexity and nature of the underlying reasoning problem.

The planned datasets include GSM8K for mathematical reasoning, CommonsenseQA for commonsense reasoning, and StrategyQA for multi-hop reasoning. A balanced subset of questions will be sampled from each dataset to ensure comparable evaluation across all selected LLMs.

\begin{table}[h]
\centering
\caption{Reasoning datasets used in the empirical evaluation.}
\label{tab:datasets}
\begin{tabular}{llll}
\hline
\textbf{Dataset} & \textbf{Task} & \textbf{Samples} & \textbf{Evaluation} \\
\hline
GSM8K & Mathematical Reasoning & 100 & Accuracy \\
CommonsenseQA & Commonsense Reasoning & 100 & Accuracy \\
StrategyQA & Multi-hop Reasoning & 100 & Accuracy \\
\hline
\end{tabular}
\end{table}

\subsection{Prompt Design}

To ensure a fair comparison across all evaluated models, a standardized prompting strategy is employed throughout the study. Every LLM receives exactly the same instruction, requiring the generation of both a final answer and a corresponding natural language explanation.

The prompt is designed to encourage concise self-generated explanations without explicitly requesting Chain-of-Thought reasoning. This decision is motivated by recent evidence suggesting that some reasoning-oriented models suppress internal reasoning traces or generate hidden reasoning processes that are not consistently accessible across different architectures. Consequently, requesting explicit Chain-of-Thought reasoning could introduce unnecessary variability between models.

The standardized prompt used in this study is shown below.

\begin{quote}
\textit{
Answer the following question. First provide your final answer. Then provide a concise explanation that justifies your answer.
}
\end{quote}

Using a single prompt template for all evaluated models minimizes prompt-induced variation and allows differences in explanation quality to be attributed primarily to model behavior rather than prompting strategy.

\subsection{Evaluation Protocol}

The proposed evaluation protocol assesses both prediction quality and explanation quality. Each generated response consists of two components: the model's final answer and its accompanying self-generated explanation. The answer is first evaluated against the ground-truth label provided by the benchmark dataset. Subsequently, the explanation is analyzed independently using a structured evaluation protocol.

The evaluation combines quantitative and qualitative assessment. Quantitative evaluation measures answer correctness using benchmark-specific evaluation metrics. Qualitative evaluation focuses on the faithfulness of the generated explanation and the identification of common explanation failure patterns.

Each explanation is manually reviewed according to a predefined faithfulness rubric to determine whether the explanation genuinely supports the predicted answer. In addition, explanations are categorized according to recurring failure patterns observed during analysis. This mixed-method evaluation enables a more comprehensive assessment of explanation quality than relying solely on prediction accuracy.

\subsection{Evaluation Metrics}

The empirical evaluation considers four complementary dimensions of explanation quality.

\begin{table}[h]
\centering
\caption{Evaluation dimensions considered in this study.}
\begin{tabular}{lp{8cm}}
\hline
Metric & Description\\
\hline
Answer Accuracy & Whether the predicted answer matches the ground-truth label.\\

Explanation Faithfulness & Whether the explanation genuinely supports the predicted answer.\\

Explanation Consistency & Stability of explanations across repeated generations for identical inputs.\\

Failure Category & Type of explanation failure identified through qualitative analysis.\\
\hline
\end{tabular}
\label{tab:metrics}
\end{table}

To evaluate explanation faithfulness, this study employs a lightweight four-level scoring rubric adapted for the purposes of this empirical analysis.

\begin{table}[h]
\centering
\caption{Faithfulness scoring rubric.}
\begin{tabular}{cl}
\hline
Score & Interpretation\\
\hline
0 & Explanation contradicts the generated answer.\\

1 & Explanation provides little or no meaningful justification.\\

2 & Explanation partially supports the predicted answer but contains missing or weak reasoning.\\

3 & Explanation clearly and consistently justifies the generated answer.\\
\hline
\end{tabular}
\label{tab:faithfulness}
\end{table}

In addition to numerical scores, qualitative analysis categorizes explanation failures into four major groups: unsupported justification, factual hallucination, incomplete reasoning, and logically inconsistent explanation. These categories provide further insight into common explanation weaknesses beyond quantitative performance measures.





\section{Experimental Setup}

\subsection{Hardware and Software Environment}

All experiments were conducted using Python 3.11 together with the Hugging Face Transformers library for model inference. Supporting libraries including PyTorch, Datasets, Pandas, NumPy, and Scikit-learn were used for data processing, evaluation, and result analysis. The experiments were executed in a Linux-based environment using NVIDIA GPU acceleration when available.

To ensure reproducibility, identical software versions and inference settings were maintained across all evaluated models. All generated responses, evaluation scores, and intermediate outputs were stored in structured CSV format for subsequent quantitative and qualitative analysis.

Table~\ref{tab:environment} summarizes the experimental environment used in this study.

\begin{table}[h]
\centering
\caption{Experimental environment.}
\label{tab:environment}
\begin{tabular}{ll}
\hline
\textbf{Component} & \textbf{Configuration} \\
\hline
Programming Language & Python 3.11 \\
Deep Learning Framework & PyTorch 2.x \\
LLM Library & Hugging Face Transformers \\
Dataset Library & Hugging Face Datasets \\
Data Processing & Pandas, NumPy \\
Evaluation & Scikit-learn \\
Operating System & Ubuntu Linux \\
Hardware & NVIDIA GPU (CUDA-enabled) \\
\hline
\end{tabular}
\end{table}

\subsection{Model Configuration}

To ensure fair comparison, all evaluated Large Language Models were executed using identical inference settings whenever supported by the respective model implementation. The same prompt template, generation parameters, and decoding strategy were applied across all experiments.

Greedy decoding was adopted to minimize response variability and improve reproducibility. Sampling-based decoding techniques such as nucleus sampling or top-k sampling were disabled, ensuring deterministic outputs across repeated runs.

The generation parameters used throughout the experiments are summarized in Table~\ref{tab:generation}.

\begin{table}[h]
\centering
\caption{Generation parameters.}
\label{tab:generation}
\begin{tabular}{ll}
\hline
\textbf{Parameter} & \textbf{Value} \\
\hline
Temperature & 0.0 \\
Top-p & 1.0 \\
Top-k & Disabled \\
Maximum New Tokens & 256 \\
Sampling & Disabled \\
Random Seed & 42 \\
Prompt Template & Standardized \\
\hline
\end{tabular}
\end{table}

\subsection{Experimental Procedure}

The experimental procedure consists of six sequential stages, illustrated conceptually in Figure~\ref{fig:pipeline}.

\begin{enumerate}
    \item Three benchmark reasoning datasets were selected and preprocessed.
    
    \item Each question was presented to every evaluated LLM using the standardized prompt.
    
    \item Each model generated both a final answer and a corresponding self-generated explanation.
    
    \item Answer correctness was automatically evaluated against the ground-truth label provided by the benchmark dataset.
    
    \item Generated explanations were assessed using the proposed evaluation protocol, including faithfulness scoring, consistency analysis, and qualitative error categorization.
    
    \item Quantitative and qualitative results were aggregated and compared across models and reasoning tasks.
\end{enumerate}

The complete experimental pipeline is shown in Figure~\ref{fig:pipeline}. The implementation has been designed to support reproducibility, enabling identical experiments to be repeated using alternative models or additional benchmark datasets.

\begin{figure}[h]
\centering
\fbox{
\parbox{0.90\linewidth}{
\centering
\textbf{Experimental Pipeline Placeholder}

\vspace{0.3cm}

Dataset Selection

$\downarrow$

Prompt Construction

$\downarrow$

LLM Inference

$\downarrow$

Answer + Self-Explanation Generation

$\downarrow$

Automatic Evaluation

$\downarrow$

Qualitative Analysis

$\downarrow$

Comparative Results
}
}
\caption{Overall experimental pipeline used in this study. The final version of the paper will replace this placeholder with a professionally designed diagram.}
\label{fig:pipeline}
\end{figure}







\section{Results}

This section presents the experimental results obtained from evaluating the selected open-source Large Language Models on multiple reasoning datasets. The evaluation combines quantitative performance measures with qualitative analysis of self-generated explanations. Quantitative results compare model performance using answer accuracy, explanation faithfulness, explanation consistency, and response characteristics, while qualitative analysis investigates common explanation failure patterns and representative examples.

\subsection{Overall Performance}

Table~\ref{tab:overall} summarizes the overall performance of all evaluated models across the three reasoning datasets. The reported metrics include answer accuracy, average explanation faithfulness, explanation consistency, and average response length. These results provide an initial comparison of the general capabilities of each model before performing a more detailed analysis.

\begin{table*}[t]
\centering
\caption{Overall performance of evaluated open-source LLMs across all reasoning datasets.}
\label{tab:overall}
\begin{tabular}{lcccc}
\hline
\textbf{Model} &
\textbf{Accuracy (\%)} &
\textbf{Faithfulness Score} &
\textbf{Consistency Score} &
\textbf{Avg. Response Length} \\
\hline
Llama 3.2 3B & -- & -- & -- & -- \\
Qwen 2.5 3B & -- & -- & -- & -- \\
Gemma 3 4B & -- & -- & -- & -- \\
DeepSeek-R1 Distill & -- & -- & -- & -- \\
\hline
\end{tabular}
\end{table*}

Figure~\ref{fig:overall_accuracy} presents a visual comparison of overall answer accuracy across all evaluated models.

\begin{figure}[h]
\centering
\fbox{\parbox{0.90\linewidth}{
\centering
Placeholder for Overall Accuracy Bar Chart
}}
\caption{Overall answer accuracy across all evaluated LLMs.}
\label{fig:overall_accuracy}
\end{figure}




\subsection{Faithfulness Evaluation}

The primary objective of this study is to evaluate whether self-generated explanations faithfully justify the corresponding model predictions. Table~\ref{tab:faithfulness_results} reports the average explanation faithfulness scores for each reasoning dataset. This analysis allows investigation of whether explanation quality varies across different reasoning domains.

\begin{table*}[t]
\centering
\caption{Average explanation faithfulness across reasoning datasets.}
\label{tab:faithfulness_results}
\begin{tabular}{lccc}
\hline
\textbf{Model} &
\textbf{GSM8K} &
\textbf{CommonsenseQA} &
\textbf{StrategyQA} \\
\hline
Llama 3.2 3B & -- & -- & -- \\
Qwen 2.5 3B & -- & -- & -- \\
Gemma 3 4B & -- & -- & -- \\
DeepSeek-R1 Distill & -- & -- & -- \\
\hline
\end{tabular}
\end{table*}

Figure~\ref{fig:faithfulness_heatmap} illustrates the average faithfulness scores across datasets using a heatmap.

\begin{figure}[h]
\centering
\fbox{\parbox{0.90\linewidth}{
\centering
Placeholder for Faithfulness Heatmap
}}
\caption{Average explanation faithfulness across reasoning datasets.}
\label{fig:faithfulness_heatmap}
\end{figure}



\subsection{Qualitative Error Analysis}

While quantitative evaluation provides an overall comparison of model performance, it does not explain why explanations fail. Therefore, a qualitative error analysis is conducted to identify recurring explanation failure patterns observed across different models and reasoning tasks.

\begin{table}[h]
\centering
\caption{Observed explanation failure categories.}
\label{tab:errors}
\begin{tabular}{lp{7cm}}
\hline
\textbf{Failure Type} & \textbf{Description}\\
\hline
Unsupported Justification & Explanation does not sufficiently support the predicted answer.\\
Hallucinated Facts & Explanation introduces information not supported by the input or general knowledge.\\
Incomplete Reasoning & Important reasoning steps are omitted or insufficiently explained.\\
Contradictory Explanation & Explanation conflicts with the predicted answer or contains internal inconsistencies.\\
\hline
\end{tabular}
\end{table}

Representative examples of explanation failures will be presented to illustrate each error category and highlight common reasoning weaknesses observed in current open-source LLMs.

\begin{figure}[h]
\centering
\fbox{\parbox{0.90\linewidth}{
\centering
Placeholder for Example Explanation Comparison
}}
\caption{Representative example comparing explanations generated by different LLMs for the same reasoning task.}
\label{fig:example}
\end{figure}

\subsection{Results Summary}

The quantitative and qualitative analyses presented in this section provide complementary perspectives on the performance of self-generated explanations. While quantitative metrics measure overall explanation quality, qualitative analysis reveals the underlying reasons for explanation failures. Together, these findings provide the foundation for the broader interpretation presented in the Discussion section.

\textit{This subsection will be updated after completing the experiments to summarize the principal findings and compare model behavior across reasoning tasks.}





\section{Discussion}

\subsection{Interpretation of Results}

The experimental results presented in the previous section provide insights into the trustworthiness of self-generated explanations produced by recent open-source Large Language Models. While answer accuracy remains an important performance indicator, the findings suggest that accurate predictions do not necessarily imply faithful explanations. Models may generate explanations that appear logically coherent and linguistically fluent while failing to accurately justify the underlying prediction.

The evaluation also highlights differences in explanation quality across reasoning tasks. Mathematical reasoning, commonsense reasoning, and multi-hop reasoning impose distinct cognitive demands on language models, which may influence both prediction accuracy and explanation faithfulness. Furthermore, qualitative analysis indicates that explanation failures often arise from incomplete reasoning, unsupported justifications, or internally inconsistent explanations rather than simple factual mistakes.

These observations reinforce the importance of evaluating explanation quality independently of prediction accuracy. Faithful explanations should be viewed as a complementary dimension of trustworthy AI rather than a by-product of strong predictive performance.

\subsection{Comparison with Previous Work}

The findings of this study are consistent with recent research suggesting that natural language explanations generated by LLMs are not always faithful representations of the underlying reasoning process. Previous studies have reported that plausible explanations can function as post-hoc rationalizations, creating the appearance of reasoning without necessarily reflecting the internal computations responsible for the final prediction.

Unlike many previous investigations that focus on individual models or proprietary systems, this work provides a unified empirical comparison across multiple recent open-source LLMs using identical datasets, prompts, and evaluation procedures. This standardized experimental design enables a more direct comparison of explanation quality while reducing methodological variation between evaluated models.

The proposed evaluation protocol also complements existing approaches by combining automatic evaluation with structured qualitative analysis. Rather than relying exclusively on prediction accuracy or human preference judgments, the study explicitly examines explanation faithfulness and recurring explanation failure patterns, providing a more comprehensive assessment of explanation quality.

\subsection{Implications for Explainable Artificial Intelligence}

The results presented in this study have several implications for Explainable Artificial Intelligence research. First, they emphasize that natural language explanations should not automatically be interpreted as evidence of faithful reasoning. Users may incorrectly assume that coherent explanations accurately describe the internal reasoning process of a language model, potentially leading to misplaced trust in AI-generated outputs.

Second, the findings suggest that evaluation of LLM explanations should extend beyond traditional performance metrics such as answer accuracy. Assessing explanation faithfulness, consistency, and common explanation failure patterns provides a more complete understanding of explanation quality and model transparency.

Finally, the study demonstrates the value of standardized empirical evaluation protocols for comparing explanation quality across different open-source LLMs. Such evaluations contribute toward the development of more trustworthy and transparent AI systems and may inform future benchmark design for explainable language models.




\section{Limitations}

Several limitations should be acknowledged when interpreting the findings of this study. First, the empirical evaluation focuses exclusively on four recent open-source instruction-tuned Large Language Models. Although these models represent diverse architectures and training strategies, the results may not generalize to proprietary commercial systems or future model generations.

Second, the evaluation is conducted using three publicly available reasoning datasets. While these benchmarks cover multiple reasoning domains, they cannot fully represent the diversity of real-world applications in which LLM explanations are used.

Third, explanation faithfulness includes a qualitative component based on a structured evaluation rubric. Although this approach enables detailed analysis of explanation quality, qualitative assessment inevitably involves a degree of subjective judgment. Future work could investigate automated or human-validated faithfulness metrics to improve evaluation reliability.

Finally, computational constraints limited the number of evaluated models and datasets. Expanding the experimental study to additional benchmark datasets, multilingual settings, and larger model families represents an important direction for future research.





\section{Future Work}

Several promising directions exist for extending this work. First, future studies may evaluate larger and more diverse collections of open-source and proprietary Large Language Models to determine whether explanation faithfulness scales with model size and reasoning capability.

Second, incorporating multimodal reasoning tasks involving both textual and visual information would provide valuable insight into explanation faithfulness beyond purely textual benchmarks.

Third, future research may investigate automated faithfulness evaluation methods that reduce reliance on manual qualitative assessment. Combining causal analysis, intervention-based evaluation, and human-centered evaluation could improve the robustness of explanation assessment.

Finally, the lightweight evaluation protocol proposed in this study could be extended into a standardized benchmark for assessing trustworthy self-generated explanations across future generations of Large Language Models.





\section{Conclusion}

This paper presented an empirical evaluation of the faithfulness of self-generated explanations produced by recent open-source Large Language Models. Motivated by growing concerns regarding the trustworthiness of natural language explanations, the study investigated whether model-generated explanations faithfully justified corresponding predictions across multiple reasoning tasks.

A comparative evaluation was proposed using multiple reasoning benchmarks, standardized prompting, and a unified evaluation protocol combining quantitative performance measures with structured qualitative analysis. Beyond measuring prediction accuracy, the study emphasized explanation faithfulness, explanation consistency, and common explanation failure patterns as important dimensions of trustworthy AI evaluation.

The findings of this work are expected to contribute to the growing body of Explainable Artificial Intelligence research by providing a reproducible evaluation methodology for self-generated explanations in open-source LLMs. More broadly, the study highlights that trustworthy AI requires not only accurate predictions but also explanations that faithfully represent the reasoning underlying model decisions.






Large Language Models produce explanations that are not always faithful
\cite{madsen2024selfexplanations}.

Chain-of-Thought prompting improves reasoning
\cite{wei2022chain}.

\bibliographystyle{unsrt}
\bibliography{references}

\end{document}


